\chapter{引言}\label{chap:introduction}

\section{研究背景与意义}

\subsection{计算显微镜：跨越时空尺度的生物学新范式}

在过去半个世纪中，结构生物学经历了从X射线晶体学到冷冻电子显微镜（Cryogenic Electron Microscopy, Cryo-EM）的革命性飞跃，使得人类能够以原子级分辨率解析单个蛋白质及其复合物的静态结构。然而，生命现象的本质在于"动态"与"相互作用"。细胞内部并非稀溶液环境，而是一个拥挤、异质且高度动态的软物质系统。大分子拥挤效应（Macromolecular Crowding）导致细胞内蛋白质的扩散行为、构象稳定性以及酶促反应动力学与体外稀溶液环境存在显著差异\cite{ellis2001macromolecular}。传统的实验手段难以在保持细胞完整性的同时，以原子分辨率捕捉毫秒级至秒级的动态过程。

在此背景下，分子动力学（Molecular Dynamics, MD）模拟作为一种"计算显微镜"（Computational Microscope），填补了实验观测在时空分辨率上的空白\cite{mccammon1977dynamics}。基于牛顿经典力学方程，MD模拟能够重现生物大分子在力场驱动下的运动轨迹。随着高性能计算（High Performance Computing, HPC）硬件——特别是图形处理器（Graphics Processing Unit, GPU）和百亿亿次（Exascale）超级计算机——的迅猛发展，以及增强采样算法和粗粒化（Coarse-Grained, CG）力场的成熟，MD模拟的研究对象已从单一蛋白质扩展至病毒衣壳（如HIV-1, SARS-CoV-2）、细胞器（如光合色素质体）乃至全细胞（Whole-Cell）尺度\cite{lenner2023whole}。

全细胞MD模拟不仅是计算生物学的终极挑战之一，更是连接分子生物学与系统生物学的桥梁。它旨在构建一个包含所有基因产物、代谢物、膜结构及溶剂环境的"数字孪生细胞"，从而在第一性原理（First Principles）层面上理解生命的涌现性（Emergence）特征\cite{feig2023current}。

\subsection{科学意义：解构生命物质的物理本质}

开展全细胞尺度的分子动力学模拟研究具有深远的科学意义和应用价值：

\textbf{阐明拥挤环境下的物理化学机制：}细胞质中的大分子浓度高达200-400 g/L，占据了约30\%的体积。这种极端的拥挤环境产生的排斥体积效应（Excluded Volume Effect）和非特异性相互作用（Quinary Structure）深刻影响着蛋白质的折叠、组装及相分离（Phase Separation）行为。全细胞模拟能够量化这些效应，揭示细胞如何利用拥挤环境调控生化反应速率\cite{spitzer2023molecular}。

\textbf{整合多组学数据并验证生物学假设：}全细胞模型的构建过程本身就是对基因组学、蛋白质组学、脂质组学及代谢组学数据的一次深度整合与自洽性验证。通过模拟，可以发现现有实验数据的缺失环节（如未知功能的基因产物），并预测生物大分子在细胞内的空间分布规律\cite{feig2023current}。

\textbf{指导合成生物学与精准药物设计：}以最小细胞（如JCVI-syn3A）为对象的模拟有助于确定维持生命所需的最小基因集，为人造细胞的设计提供蓝图。同时，全病毒或全细胞模拟能够揭示药物分子在复杂生物环境中的输运路径及与靶点的动态结合机制（如HIV-1衣壳抑制剂Lenacapavir的作用机理），为开发新型抗病毒药物提供理论依据\cite{lenner2023whole}。

\textbf{推动计算科学的发展：}全细胞模拟涉及数亿至数十亿粒子的计算，对超级计算机的并行效率、数据存储及可视化技术提出了极限挑战，将直接推动Exascale计算技术和大数据分析算法的革新\cite{feig2023current}。

\subsection{领域现状：通往全细胞模拟的里程碑}

近年来，该领域取得了一系列标志性成果，为全细胞模拟奠定了坚实基础：

\textbf{病毒衣壳的全原子模拟：}2013年，UIUC团队实现了HIV-1成熟衣壳（约6400万原子）的全原子模拟，揭示了衣壳的力学性质和离子通透性\cite{zhao2013mature}。随后，SARS-CoV-2病毒包膜（约3.05亿原子）的模拟揭示了刺突蛋白的糖基化屏蔽机制和动态构象\cite{casalino2020all}。

\textbf{细菌细胞质的拥挤模拟：}RIKEN和密歇根州立大学利用GENESIS软件实现了大肠杆菌细胞质的大规模全原子模拟，发现代谢物与蛋白质的相互作用会改变蛋白质的表面扩散行为\cite{feig2023current}。

\textbf{最小细胞JCVI-syn3A的初步尝试：}J. Craig Venter研究所（J. Craig Venter Institute, JCVI）与UIUC合作，构建了JCVI-syn3A的马蒂尼（Martini）粗粒化模型（约6亿粒子），并结合反应-扩散主方程（Reaction-Diffusion Master Equation, RDME）实现了首个包含完整代谢和遗传信息处理过程的全细胞动力学模型\cite{lenner2023whole}。

\section{长程静电求解的挑战与机遇}

\subsection{传统方法的通信瓶颈}

在分子动力学模拟中，长程静电相互作用的计算是最为关键的性能瓶颈之一。传统的Ewald求和方法及其网格加速版本——粒子网格Ewald（Particle Mesh Ewald, PME）\cite{darden1993particle}和粒子-粒子粒子-网格（Particle-Particle Particle-Mesh, PPPM）\cite{hockney1988computer}——通过将长程相互作用分解为短程和长程两个部分，实现了$O(N\log N)$的计算复杂度。然而，这些方法依赖于全局快速傅里叶变换（Fast Fourier Transform, FFT），在大规模并行计算中面临严重的通信瓶颈。

具体而言，FFT要求在每个时间步进行全局数据交换，这在分布式内存架构的超级计算机上会产生显著的通信开销。随着体系规模的增大和处理器数量的增加，通信时间可能超过计算时间，导致并行效率急剧下降。这一问题在全细胞模拟场景中尤为突出，因为全细胞体系的粒子数可能达到数十亿量级\cite{feig2023current}。

\subsection{随机批量方法的兴起}

近年来，随机批量方法（Random Batch Methods, RBM）为大规模粒子系统模拟提供了新的思路。Jin等人提出的随机批量相互作用粒子系统方法\cite{jin2020random}，通过随机采样策略将$O(N^2)$的相互作用计算降低为$O(N)$，同时保持了统计无偏性。这一方法已被成功应用于分子动力学、朗之万动力学和Vlasov方程等领域。

随机批量方法的核心思想是：在每个时间步，随机选择一部分相互作用进行计算，而非计算所有相互作用。通过精心设计的采样策略和适当的统计校正，可以在保持数值稳定性的前提下显著降低计算成本。更重要的是，随机批量方法天然支持并行化，因为它避免了全局通信的需求\cite{jin2020random}。

\subsection{高斯和分解的数值优势}

高斯和分解（Sum-of-Gaussians, SOG）是一种将复杂函数表示为多个高斯函数加权和的技术。在量子力学中，SOG张量神经网络已被成功应用于高维薛定谔方程的求解\cite{guo2025sum}。SOG分解的优势在于：

（1）高斯函数具有良好的解析性质，便于数值计算；
（2）高斯函数的可分离性使得高维积分可以分解为低维积分的乘积；
（3）高斯函数的快速衰减性保证了数值稳定性；
（4）SOG分解可以与随机采样策略自然结合。

将SOG分解应用于长程静电相互作用的计算，有望实现更平滑的数值行为和更好的并行扩展性。

\section{研究目标与内容}

\subsection{核心研究目标}

本研究的核心目标是探索一种兼顾数值平滑性、统计稳定性与算法扩展潜力的新型长程相互作用处理框架，并以基于随机批量采样与高斯和分解的长程静电方法（Random Batch Sum-of-Gaussians, RBSOG）作为当前阶段的核心研究对象。具体目标包括：

（1）构建面向膜体系的分子动力学模拟原型系统，实现RBSOG方法的原型验证；
（2）建立系统化的基准测试体系，从计算性能、统计方差和结构稳定性三个维度评估RBSOG方法；
（3）通过实验分析揭示RBSOG方法的优势与局限，为后续工程优化提供依据；
（4）建立可复现的自动化实验与论文生成流程，为后续研究奠定基础。

\subsection{论文组织结构}

本文的组织结构如下：

第一章为引言，介绍研究背景、科学意义、领域现状以及长程静电求解的挑战与机遇。

第二章为方法，详细介绍分子动力学基础理论、RBSOG算法框架、膜体系构建方法以及评测指标设计。

第三章为结果，展示PPPM与RBSOG的基准测试结果、批量大小扫描实验、结构稳定性分析以及JIT消融实验。

第四章为讨论，对实验结果进行深入分析，讨论RBSOG方法的优势与局限，以及与现有工作的比较。

第五章为结论，总结本研究的主要贡献和创新点，展望未来研究方向。
